\section{Source Section III.F: Boiling in Gravity}
\label{sec:boiling-gravity}

\sourceeqrange{Eq. (3.20), figures 11--13}

\subsection*{Purpose}

This subsection records the gravity parameters needed to read the boiling
examples, while keeping those finite estimates separate from the computed bubble
dynamics.

\subsection*{Reading Guide}

The gravity parameter \(g^*=mg\ell/\epsilon\) is introduced in Section III.A and
appears in the scaled total-energy equation in Appendix B.  The boiling scenario
uses a nonzero gravity branch together with the bottom/top temperature boundary
branch.  Equation (3.20) gives a bulk adiabatic temperature-gradient estimate;
in this companion it is used only as a dimensional and interpretive check.

The two finite quantities carried into this reading guide are
\[
\begin{array}{c|c|c}
\text{source anchor} & \text{quantity} & \text{typed role} \\
\hline
\text{Eq. (3.7)} & g^*=mg\ell/\epsilon
  & \text{dimensionless body-force branch parameter} \\
\text{Eq. (3.20)} &
  \rho g\left(\partial T/\partial p\right)_s
  & \text{adiabatic temperature-gradient estimate}
\end{array}
\]
The first quantity enters the scaled balance equations.  The second tells the
reader what units and sign information to expect from a bulk adiabatic-gradient
estimate.  Neither one determines a boiling onset criterion by itself.

\subsection*{Finite Checks}

The checkable statements are that \(g^*\) is dimensionless and that the product
\(\rho g(\partial T/\partial p)_s\) has the units of a temperature gradient.
The companion does not compute bubble periods, Reynolds numbers, or heat-flux
peaks.

The adiabatic-gradient estimate starts with hydrostatic pressure balance in the
vertical coordinate \(z\):
\begin{equation}
  \frac{\dd p}{\dd z}=-\rho g.
\end{equation}
Along an adiabatic bulk displacement,
\begin{equation}
  \frac{\dd T}{\dd z}
  =\left(\frac{\partial T}{\partial p}\right)_s
   \frac{\dd p}{\dd z},
\end{equation}
so
\begin{equation}
  \left(\frac{\dd T}{\dd z}\right)_{\rm bulk}
  =-\rho g\left(\frac{\partial T}{\partial p}\right)_s.
  \sourceeqbadge{3.20}
  \label{eq:iii-adiabatic-gradient-derived}
\end{equation}
Here \(\rho g\) has pressure-gradient units and
\((\partial T/\partial p)_s\) converts pressure to temperature at fixed entropy.
The product is therefore a temperature gradient.

For the order estimate printed with Eq. (3.20), use \(\rho=mn\),
\(g^*=mg\ell/\epsilon\), a pressure scale \(p_c\sim\epsilon/v_0\), and
\(n\sim1/v_0\).  Then
\begin{equation}
  \rho g\left(\frac{\partial T}{\partial p}\right)_s
  \sim \frac{mngT_c}{p_c}
  \sim \frac{mg\ell}{\epsilon}\frac{T_c}{\ell}
  =g^*\frac{T_c}{\ell}.
\end{equation}
Thus the source order statement is a scale estimate, not an exact equality.
It does not determine a boiling threshold or a bubble trajectory.  Its role is
to give the reader a bulk scale against which source temperature fields can be
interpreted when the gravity branch is active.
The gravity and adiabatic-gradient dimensional checks are covered by the
Section III mirror pair listed in Section~\ref{sec:numerical-results-overview}.
They are not boiling onset calculations.

\subsection*{Limitations}

Boiling onset, transition between weak and strong boiling, and fully developed
boiling are simulation observations.  They are outside the verification scope of
this guide.
